% ============================================================
% INFORME DE IMPACTO SOCIAL, AMBIENTAL Y TERRITORIAL
% Proyecto: Pista de Atletismo — Isla de Pascua (Rapa Nui)
% Investigadores: Fabián Ortega, Lucas Nervi, Vicente Belgeri & Matías Pardo (2026)
% ============================================================

\documentclass[12pt, a4paper]{article}

% --- Paquetes básicos ---
\usepackage[utf8]{inputenc}
\usepackage[T1]{fontenc}
\usepackage[spanish, es-noshorthands]{babel}
\usepackage{lmodern}
\usepackage{microtype}

% --- Márgenes y geometría ---
\usepackage[top=2.5cm, bottom=2.5cm, left=3cm, right=2.5cm, headheight=14pt]{geometry}

% --- Tipografía ---
\usepackage{charter}   % Fuente serif legible para documentos institucionales

% --- Encabezados y pies de página ---
\usepackage{fancyhdr}
\pagestyle{fancy}
\fancyhf{}
\rhead{\small Informe de Impacto — Pista de Atletismo Rapa Nui}
\lhead{\small Ortega, Nervi, Belgeri \& Pardo, 2026}
\cfoot{\thepage}
\renewcommand{\headrulewidth}{0.4pt}

% --- Listas y tablas ---
\usepackage{enumitem}
\usepackage{booktabs}
\usepackage{array}
\usepackage{tabularx}

% --- Colores e hipervínculos ---
\usepackage[dvipsnames]{xcolor}
\usepackage[colorlinks=true, linkcolor=NavyBlue, citecolor=NavyBlue,
            urlcolor=NavyBlue, bookmarks=true]{hyperref}

% --- Gráficos y cajas ---
\usepackage{graphicx}
\usepackage{tcolorbox}
\tcbuselibrary{skins, breakable}

% --- Cajas de citas etnográficas ---
\newtcolorbox{citaetno}{
  breakable,
  enhanced,
  colback=gray!8,
  colframe=gray!40,
  boxrule=0.5pt,
  left=10pt, right=10pt, top=6pt, bottom=6pt,
  fontupper=\itshape,
  before upper={\footnotesize\noindent\rule{\linewidth}{0.3pt}\\[4pt]},
  after upper={\\[4pt]\rule{\linewidth}{0.3pt}}
}

% --- Caja de concepto clave ---
\newtcolorbox{conceptoclave}[1][]{
  breakable,
  enhanced,
  colback=NavyBlue!5,
  colframe=NavyBlue!50,
  boxrule=0.8pt,
  title={Concepto clave},
  fonttitle=\bfseries\small,
  #1
}

% --- Bibliografía ---
\usepackage[authoryear, round, sort]{natbib}
\bibliographystyle{apalike}

% --- Párrafos y espaciado ---
\setlength{\parindent}{0pt}
\setlength{\parskip}{8pt}
\renewcommand{\baselinestretch}{1.2}

% --- Sección personalizada ---
\usepackage{titlesec}
\titleformat{\section}{\large\bfseries\color{NavyBlue}}{\thesection.}{0.6em}{}[\titlerule]
\titleformat{\subsection}{\normalsize\bfseries}{\thesubsection.}{0.5em}{}
\titleformat{\subsubsection}{\normalsize\itshape}{\thesubsubsection.}{0.4em}{}

% ============================================================
%                    INICIO DEL DOCUMENTO
% ============================================================
\begin{document}

% ---- Portada ----
\begin{titlepage}
  \centering
  \vspace*{2cm}
  {\Large\bfseries INFORME DE IMPACTO SOCIAL, AMBIENTAL Y TERRITORIAL\par}
  \vspace{0.5cm}
  {\rule{\linewidth}{1.5pt}\par}
  \vspace{0.8cm}
  {\LARGE\bfseries\color{NavyBlue}
    Proyecto: Pista de Atletismo\\[0.3cm]
    Isla de Pascua (Rapa Nui)\par}
  \vspace{1.5cm}
  {\large Informe Preliminar — Sistema Nacional de Inversiones (SNI)\par}
  \vspace{2cm}

  \begin{minipage}{0.9\textwidth}
    \centering
    {\normalsize\textbf{Elaborado por:}\\[0.4cm]
    Fabián Ortega \& Lucas Nervi\\[0.1cm]
    \textit{En base a investigación cualitativa de campo (2024--2026)}\\[0.1cm]
    {\footnotesize Metodología: Teoría Fundamentada Constructivista con principios polinésicos (Talanoa / Vā)}\\[0.4cm]
    \rule{0.5\linewidth}{0.3pt}\\[0.3cm]
    Vicente Belgeri \& Matías Pardo\\[0.1cm]
    {\footnotesize Perfil estratégico, justificación institucional y formulación SNI}\par}
  \end{minipage}

  \vspace{2cm}
  {\normalsize Hanga Roa, Isla de Pascua\par}
  {\normalsize Marzo 2026\par}

  \vfill
  {\footnotesize\color{gray}
    Este informe se fundamenta en una investigación cualitativa con 17 actores clave
    del ecosistema deportivo y comunitario de Rapa Nui (2024--2025).
    Los hallazgos y citas están sujetos a validación comunitaria.\par}
\end{titlepage}

% ---- Tabla de contenidos ----
\tableofcontents
\newpage

% ============================================================
\section{Introducción}
% ============================================================

La comuna de Isla de Pascua (Rapa Nui) constituye uno de los territorios habitados
permanentemente más aislados del planeta, ubicada a aproximadamente 3.700 km del
territorio continental de Chile. Con una superficie de 163,6 km$^2$ y una población
cercana a 7.750 habitantes ---concentrada principalmente en el poblado de Hanga Roa---,
Rapa Nui representa un caso singular dentro del sistema de inversión pública chilena:
un territorio insular con reconocimiento de pueblo indígena originario y una historia
colonial documentada que sigue determinando sus condiciones estructurales actuales.

La economía local depende principalmente del turismo, la pesca artesanal y las
actividades culturales asociadas al patrimonio rapanui. En este contexto, la
infraestructura deportiva cumple un rol estratégico de múltiples dimensiones que
trasciende la provisión de un servicio: actúa simultáneamente como dispositivo de
salud pública, de organización comunitaria, de transmisión cultural y de resiliencia
identitaria \citep{hauofa1994}.

El presente informe tiene como objetivo evaluar el \textbf{impacto social, ambiental
y territorial de la construcción de una pista de atletismo en Rapa Nui}, con el fin
de justificar su pertinencia dentro del marco de inversión pública. A diferencia de
informes estándar de este tipo, este documento se fundamenta en una investigación
cualitativa de campo con 17 actores clave del ecosistema deportivo y comunitario de
la isla, desarrollada durante 2024--2025 mediante \textbf{Teoría Fundamentada
Constructivista} \citep{charmaz2014} y principios ético-metodológicos polinésicos
(\textbf{Talanoa} y \textbf{Vā}) \citep{prescott2008, kaili2005}.

\subsection{Nota metodológica}

Las afirmaciones sobre impacto social de este informe se derivan de evidencia
empírica recolectada en terreno, no de supuestos genéricos. Las citas textuales
que aparecen a lo largo del documento provienen de entrevistas en profundidad con
deportistas, entrenadores, dirigentes, familias y profesionales de salud y educación
de Rapa Nui. Los principios de diseño derivados del análisis están sujetos a
\textbf{validación comunitaria} antes de cualquier implementación.

% ============================================================
\section{Descripción del Proyecto}
% ============================================================

El proyecto consiste en la \textbf{construcción de una pista de atletismo de estándar
competitivo}, orientada a fomentar la práctica deportiva, el desarrollo del atletismo
y la realización de actividades comunitarias en la comuna de Isla de Pascua.

\subsection{Componentes principales}

\begin{itemize}[leftmargin=*, itemsep=3pt]
  \item Pista de atletismo de 400 metros con superficie sintética
  \item Áreas para disciplinas atléticas (saltos y lanzamientos)
  \item Iluminación deportiva
  \item Graderías y áreas de espectadores
  \item Equipamiento deportivo especializado
  \item Espacios de entrenamiento comunitario y multiuso
\end{itemize}

\subsection{Objetivos del proyecto}

\begin{itemize}[leftmargin=*, itemsep=3pt]
  \item Fomentar la práctica deportiva en niños, jóvenes y adultos de la comunidad
  \item Generar infraestructura para el desarrollo del atletismo competitivo y otras
        disciplinas
  \item Promover hábitos de vida activa y saludable en la comunidad rapanui
  \item Crear un espacio comunitario de encuentro social y cultural pertinente al
        territorio
  \item Apoyar el desarrollo de talentos deportivos locales con condiciones mínimas
        dignas
\end{itemize}

% ============================================================
\section{Diagnóstico Territorial}
\label{sec:diagnostico}
% ============================================================

\subsection{Marco histórico-colonial}

Para comprender el impacto de cualquier inversión pública en Rapa Nui, es
indispensable situar el territorio en su historia. La experiencia contemporánea del
pueblo rapanui está moldeada por una secuencia de traumas históricos que
desmantelaron la sociedad pre-contacto e instalaron un sistema colonial de
explotación y control social \citep{cristino2011}.

Entre 1862 y 1863, redadas esclavistas peruanas se llevaron a la fuerza aproximadamente
al 35\% de la población. El retorno de supervivientes introdujo epidemias que exterminaron
a más del 90\% de la población restante. Para 1877, solo sobrevivían 111 personas en
la isla. El período de la Compañía Explotadora de Isla de Pascua (CEIP, 1895--1953)
consolidó un régimen de \textbf{confinamiento forzado en Hanga Roa}, despojando a la
comunidad rapanui de sus tierras y reduciéndola a condición de servidumbre.

Esta historia no es meramente un telón de fondo: es \textbf{una memoria activa y
viva} que informa la identidad, la política y las relaciones comunitarias actuales,
incluyendo la relación de la comunidad con las instituciones del Estado chileno
\citep{foerster2021}. El conflicto fundamental sobre la interpretación del ``Acuerdo
de Voluntades'' de 1888 ---que para la comunidad rapanui fue un tratado de protección
que reservaba sus derechos sobre el territorio--- permanece \textbf{irresoluto} hasta
el presente.

\subsection{Aislamiento geográfico y condiciones logísticas}

La insularidad impone condiciones estructurales que no tienen equivalente en el
territorio continental. La investigación cualitativa que sustenta este informe
identificó las siguientes condicionantes logístico-económicas \citep{ortega2026}:

\begin{itemize}[leftmargin=*, itemsep=4pt]
  \item Altos costos de traslado para competencias: cada atleta que sale a competir
        al continente requiere pasajes aéreos de aproximadamente \$700.000--\$1.200.000
        sólo en transporte
  \item Imposibilidad de cumplir calendarios de competencia diseñados para el
        continente: federaciones nacionales exigen participación en torneos que
        logísticamente son inviables desde la isla
  \item Infraestructura deportiva limitada y envejecida, sin mantención adecuada
  \item Dependencia de recursos fluctuantes y apoyos personales intermitentes
        para el funcionamiento deportivo
\end{itemize}

\subsection{Demografía y actores del deporte}

\begin{itemize}[leftmargin=*, itemsep=3pt]
  \item Población aproximada: 7.750 habitantes (Censo 2017)
  \item Mayoría concentrada en Hanga Roa
  \item Presencia significativa de población rapanui con reconocimiento como
        pueblo indígena originario
  \item Ecosistema deportivo activo pero fragmentado: múltiples disciplinas
        (atletismo, natación, fútbol, ciclismo, deportes tradicionales) sin
        coordinación institucional eficaz
\end{itemize}

\subsection{Fragmentación institucional y brechas actuales}

La investigación cualitativa evidenció que la oferta deportiva en Rapa Nui es
\textbf{fragmentada}: cada disciplina negocia por separado con federaciones nacionales,
la gestión es débil y voluntarista, y hay escasez virtual de apoyo multidisciplinario
(prevención de lesiones, salud mental deportiva, nutrición) \citep{ortega2026}.

Esta fragmentación tiene consecuencias concretas: agotamiento de dirigentes
voluntarios que enfrentan burocracia diseñada para organizaciones profesionalizadas,
dependencia de liderazgos carismáticos individuales que cuando migran o se retiran
colapsan las iniciativas, y refuerzo de la desconfianza histórica hacia instituciones
externas.

\begin{conceptoclave}
  \textbf{Régimen colonial-burocrático} (concepto emergente de la investigación):
  Las normativas deportivas nacionales, los costos de competencia continental y las
  lógicas organizativas diseñadas para el ``continente'' operan como barrera
  estructural para el desarrollo deportivo insular. No es simplemente una cuestión
  de recursos limitados, sino de \textit{asimetría epistémica}: las reglas del juego
  fueron escritas sin considerar la realidad insular \citep{escobar2007}.
\end{conceptoclave}

% ============================================================
\section{Identificación del Problema Público}
% ============================================================

\subsection{Problema central}

La comuna de Rapa Nui presenta un \textbf{déficit de infraestructura deportiva
especializada para la práctica del atletismo}, lo que limita el acceso regular,
seguro y de calidad a una disciplina base para el desarrollo físico, educativo
y comunitario de la población. Este déficit no puede analizarse de forma aislada:
se inscribe en una historia de postergación institucional y aislamiento geográfico
que refuerza brechas estructurales documentadas por la investigación cualitativa
\citep{ortega2026}.

\subsection{Manifestaciones del problema}

El problema descrito se expresa en diversas situaciones observables:

\begin{itemize}[leftmargin=*, itemsep=3pt]
  \item Ausencia de una pista reglamentaria o semirreglamentaria para entrenamiento atlético
  \item Uso de espacios no diseñados para pruebas atléticas, con riesgo de lesiones
  \item Limitaciones severas para actividades escolares de atletismo (educación física, torneos)
  \item Menor proyección de talentos deportivos locales que deben migrar al continente
  \item Baja capacidad para organizar competencias, encuentros o concentraciones
  \item Pérdida de oportunidades para vincular deporte con turismo, salud y comunidad
\end{itemize}

\subsection{Causas estructurales}

Entre las principales causas se identifican:

\begin{itemize}[leftmargin=*, itemsep=3pt]
  \item Rezago histórico en infraestructura deportiva especializada, agravado por la condición colonial
  \item Alta complejidad logística y sobrecosto de materiales en territorio insular
  \item Priorización de otras necesidades públicas urgentes
  \item Fragmentación institucional del ecosistema deportivo local (documentada en Sección~\ref{sec:diagnostico})
  \item Ausencia de un relato estratégico que conecte la pista con beneficios multisectoriales
\end{itemize}

\subsection{Efectos sociales del problema}

La carencia de infraestructura produce efectos de corto y largo plazo:

\begin{itemize}[leftmargin=*, itemsep=3pt]
  \item Menor acceso a actividad física estructurada y de calidad
  \item Menor adherencia deportiva de niños y jóvenes (deserción documentada)
  \item Pérdida de oportunidades de prevención en salud física y mental
  \item Debilitamiento de trayectorias de formación deportiva sostenida
  \item Subutilización del deporte como herramienta de integración y cohesión social
\end{itemize}

% ============================================================
\section{Impacto Social}
% ============================================================

\subsection{Salud pública}

La investigación cualitativa reveló que el \textbf{deporte como medio para mejorar
la salud física} emerge como una narrativa central y proactiva en la comunidad rapanui.
Sin embargo, esta narrativa choca con brechas estructurales: déficit de kinesiólogos
y apoyo multidisciplinario obliga a la comunidad a aprender ``a porrazo'' por dolores
\citep{ortega2026}.

El análisis de la categoría Salud identificó un perfil epidemiológico preocupante en
la isla: alta prevalencia de enfermedades crónicas vinculadas a cambios alimentarios
post-colonización, y una crisis de salud mental documentada con 560 consultas en el
Hospital Hanga Roa en 2018 ---la cifra más alta, superando a hipertensión y diabetes.
La infraestructura deportiva de calidad puede actuar como dispositivo preventivo en
ambas dimensiones.

El deporte promueve además un cambio cultural específico documentado en las entrevistas:
\textbf{el cuidado de lesiones deja de verse como debilidad} y el descanso como
parte del entrenamiento es revalorizado, contribuyendo a una mayor conciencia corporal
en la comunidad.

\subsection{Formación deportiva y trayectorias}

La investigación identificó que la variación en trayectorias deportivas ---por qué
algunas son sostenidas y otras intermitentes o truncas--- \textbf{no depende de la
motivación individual}, sino de la capacidad colectiva de la comunidad para activar
condiciones estructurales.

Una pista de atletismo de estándar competitivo habilitaría concretamente:

\begin{itemize}[leftmargin=*, itemsep=4pt]
  \item \textbf{Fase de iniciación (6--12 años):} espacio lúdico multiactividad sin
        especialización temprana, reduciendo la brecha de falta de oferta competitiva
        infantil documentada
  \item \textbf{Fase de desarrollo (13--17 años):} introducción gradual a la
        especificidad sin necesidad de traslado al continente
  \item \textbf{Fase de especialización (18+):} condiciones mínimas para atletas
        locales que actualmente entrenan en condiciones precarias
\end{itemize}

La ausencia de apoyos sostenidos genera deserción. Un entrevistado relató:

\begin{citaetno}
``Éramos 4 competidores, yo tuve que invitar a mi hermano que ya no competía,
a mi otra hermana que no competía, yo y había otros chicos más y los chicos
los que estaban.''\\[4pt]
\rightline{\footnotesize--- Atleta participante, entrevista en profundidad}
\end{citaetno}

Esta cita ilustra cómo la falta de infraestructura y estímulos suficientes provoca
baja participación, obligando a recurrir a personas que ya habían abandonado el
deporte.

\subsection{Cohesión comunitaria y capital social}

La investigación documentó que la práctica deportiva actúa como \textbf{motor de
cohesión comunitaria} en Rapa Nui. El ciclo virtuoso se alimenta de la
\textit{persistencia en la participación}, que genera \textit{percepción de
competencia y logro}, reforzando el \textit{valor de la participación} \citep{ortega2026}.

El análisis del \textbf{capital social y estructuras de apoyo} identificó que cuando
los mecanismos de apoyo colectivo operan efectivamente, se observan:

\begin{itemize}[leftmargin=*, itemsep=3pt]
  \item Mayor retención de atletas y participantes
  \item Regreso de talentos con aprendizaje externo (migración circular no definitiva)
  \item Aumento de participación femenina
  \item Reputación positiva de las organizaciones deportivas en la comunidad
\end{itemize}

\subsection{Oportunidades para la juventud}

El discurso sobre la juventud en Rapa Nui se articula en torno a una poderosa
narrativa de potencialidad. Los jóvenes son concebidos como \textbf{``Diamantes en
Bruto''}, portadores de una energía juvenil como recurso fundamental para el
desarrollo colectivo \citep{ortega2026}. Sin embargo, esta narrativa coexiste con
\textbf{tensiones intergeneracionales}:

\begin{citaetno}
``Eso es pa' nosotros, ya pasó, pa' ustedes ya tiene todo esto, tiene beca tiene todo.''\\[4pt]
\rightline{\footnotesize--- Madre participante, entrevista en profundidad}
\end{citaetno}

Esta frase ilustra ---desde una perspectiva, no la única--- cómo el deporte puede
ser visto por algunas familias como ``superado'' frente a la educación formal. Sin
embargo, la mayoría de los actores entrevistados reconocen que \textbf{educación y
deporte pueden y deben coexistir}: el acceso a infraestructura deportiva de calidad
contribuye a reducir conductas de riesgo, fomenta la disciplina y el trabajo colectivo,
y genera oportunidades de desarrollo que complementan la formación académica.

\subsection{Motivación y adherencia deportiva}

El análisis de la categoría \textbf{Motivación Deportiva} reveló que la motivación
en Rapa Nui se entiende como un equilibrio delicado entre la fuerza que nace de la
persona misma y los apoyos que provienen del entorno. La disciplina, la constancia
y la automotivación grupal son factores decisivos, pero no suficientes:

\begin{itemize}[leftmargin=*, itemsep=3pt]
  \item La motivación personal existe y es sólida, pero se ve condicionada por
        condiciones materiales: sin infraestructura, sin incentivos, sin apoyo,
        la disciplina individual no alcanza para sostener trayectorias
  \item Las \textit{aspiraciones deportivas} trascienden lo local: deseo de
        proyección nacional e internacional funciona como horizonte simbólico
  \item Sin condiciones externas adecuadas, emerge la \textbf{deserción deportiva
        por carencia de apoyo y logros} o la transición de atleta a emprendedor
        por necesidad económica
\end{itemize}

% ============================================================
\section{Impacto en Identidad y Cultura Rapanui}
% ============================================================

Esta sección constituye el aporte más diferenciador de este informe respecto a
estudios estándar de infraestructura deportiva. La cultura rapanui no es un
``factor de riesgo a mitigar'': es \textbf{el contexto que da sentido y
sostenibilidad} a cualquier proyecto de esta naturaleza.

\subsection{Hanai y deporte como práctica comunitaria}

El concepto rapanui de \textbf{\textit{hanai}} representa una concepción no biológica
de la crianza: la responsabilidad del cuidado y la formación de los jóvenes es un
\textit{compromiso social colectivo}, no un deber exclusivo de los padres biológicos.
Tíos, abuelos y otros miembros de la familia extendida participan activamente en la
formación de los jóvenes \citep{pakomio2016}.

Este principio se traduce directamente en la práctica deportiva: \textbf{la comunidad
concibe el desarrollo del atleta joven como responsabilidad colectiva}. Una pista de
atletismo no es solo equipamiento: es un espacio donde este compromiso comunitario
se materializa físicamente.

\subsection{Los \textit{koros} como transmisores de saber deportivo}

La investigación documentó que los \textbf{\textit{koros}} (abuelos, ancianos
conocedores) no son meros observadores del deporte: son \textbf{portadores de saberes
técnicos deportivos} equivalentes al entrenamiento formal, transmitidos
intergeneracionalmente a través de práctica y relato \citep{ortega2026}.

Esta dualidad de saberes refleja la construcción de un \textit{habitus} deportivo
híbrido \citep{bourdieu1990}: el cuerpo aprende tanto a través del entrenamiento
técnico-científico como de la autogestión y la disciplina comunitaria. Una
infraestructura que reconoce este saber ---permitiendo que los \textit{koros} sean
parte del ecosistema de entrenamiento--- fortalece la pertinencia cultural del
proyecto.

\subsection{Tapati Rapa Nui y el deporte como acto identitario}

El \textbf{Tapati Rapa Nui} es el festival cultural anual que incluye importantes
componentes deportivos y de destreza física como el \textit{Haka Pei} ---un rito de
valentía que consiste en deslizarse por la ladera de un cerro en un trineo de
troncos de plátano. Esta práctica, junto con la memoria del \textit{Tangata Manu}
(Hombre Pájaro), documenta que el deporte en Rapa Nui \textbf{es un acto identitario
antes que una práctica meramente recreativa} \citep{ortega2026, cristino2011}.

La infraestructura deportiva moderna, para ser pertinente, debe dialogar con esta
tradición: no reemplazarla, sino crear condiciones para que las prácticas deportivas
actuales también sean vividas como expresiones de identidad rapanui.

\subsection{El espacio relacional (\textit{Vā}) y la infraestructura}

El concepto polinésico de \textbf{\textit{Vā}} refiere al espacio relacional entre
personas que debe ser cuidado y nutrido \citep{kaili2005}. Aplicado a la
infraestructura: una pista de atletismo no es solo un espacio físico, sino un
\textbf{espacio relacional} donde se construyen y fortalecen vínculos comunitarios.

Esto tiene implicaciones concretas para el diseño:

\begin{itemize}[leftmargin=*, itemsep=3pt]
  \item Los espacios de graderías y áreas comunitarias son tan importantes como
        la pista misma: son donde el \textit{Vā} se activa
  \item El proceso de diseño debe ser participativo: la comunidad debe poder
        ``hacer propio'' el espacio antes de que exista físicamente
  \item Los usos del espacio deben incluir actividades culturales, no solo
        competencias deportivas técnicas
\end{itemize}

\subsection{Mar de islas: identidad y proyección deportiva}

El filósofo polinésico Epeli Hau'ofa \citeyearpar{hauofa1994} propone que las islas
del Pacífico no son pequeños territorios aislados, sino nodos de un vasto ``mar de
islas'': la identidad isleña es fluida, conectada y resiliente. El deportista rapanui
se mueve en este océano de posibilidades, \textbf{conectando su isla con otros mundos
sin perder su centro}.

Una pista de atletismo de estándar competitivo ampliaría este horizonte: permitiría
a atletas rapanui prepararse en la isla y proyectarse al continente y a la región del
Pacífico sin necesidad de abandonar definitivamente su comunidad.

% ============================================================
\section{Justificación Estratégica del Proyecto}
% ============================================================

\subsection{Más que una pista: infraestructura estructurante}

La principal fortaleza de este proyecto es que no debe ser comunicado como una
obra aislada, sino como una \textbf{infraestructura estructurante} para el
desarrollo social, deportivo y territorial de Rapa Nui. La investigación
cualitativa que sustenta este informe lo confirma: la comunidad no concibe
la pista como un equipamiento técnico, sino como un \textit{espacio de
posibilidad colectiva} que habilita cinco mecanismos simultáneamente:
legitimidad, profesionalización, redes, cuidado integral y gobernanza
\citep{ortega2026}.

La narrativa recomendada para la postulación ante organismos público es:

\begin{quote}
\textit{``La pista de atletismo de Rapa Nui no sólo responde a una necesidad
deportiva. Se trata de una infraestructura pública habilitante para la salud,
la formación de jóvenes, la cohesión comunitaria y la proyección del
territorio como referente insular de deporte, bienestar y encuentro en el
Pacífico Sur.''}
\end{quote}

\subsection{Equidad territorial e inversión con retornos múltiples}

El acceso a infraestructura pública de calidad no debe depender exclusivamente
de la cercanía con grandes centros urbanos. En territorios insulares y alejados,
cada obra pública tiene un efecto redistributivo más intenso: corrige brechas
históricas en acceso a equipamiento, servicios y oportunidades de desarrollo personal.

El proyecto posee retornos de naturaleza múltiple:

\begin{itemize}[leftmargin=*, itemsep=3pt]
  \item \textbf{Retorno social:} mejora de bienestar, salud y convivencia
  \item \textbf{Retorno educativo:} fortalecimiento de trayectorias escolares y deportivas
  \item \textbf{Retorno económico:} eventos, turismo de escala compatible y servicios asociados
  \item \textbf{Retorno simbólico:} señal potente de inversión en un territorio de alta singularidad
  \item \textbf{Retorno político-institucional:} proyecto visible, defendible y alineado con equidad territorial
\end{itemize}

\subsection{Proyección estratégica: Centro Atlético del Pacífico Sur}
\label{sec:centro_pacifico}

La formulación recomendada no es \textit{``solo una pista''}. El enfoque
estratégico consiste en presentar la obra como la primera infraestructura
habilitante de un \textbf{Centro Atlético del Pacífico Sur}, concebido como
polo de formación atlética, intercambio deportivo, deporte escolar insular,
bienestar y vida activa, e investigación aplicada a deporte en contextos remotos.

Este cambio de escala mejora significativamente el peso político del proyecto
porque conecta la obra con una visión de futuro, aumenta la capacidad de atraer
apoyos interinstitucionales y permite construir un relato diferenciador.
A mediano plazo, el Centro Atlético del Pacífico Sur podría incorporar sala
multipropósito para preparación física, unidad de evaluación deportiva, espacios
de educación y clínica deportiva, residencias cortas para concentraciones, y
colaboración con universidades y programas de salud.

\subsection{Ventaja habilitante: emplazamiento en Fundo Vaitea}

Un elemento especialmente favorable para la maduración de esta iniciativa es la
existencia de un terreno propuesto en el sector de \textbf{Fundo Vaitea}, sobre
terrenos donados para efectos del proyecto. La disponibilidad preliminar de suelo
representa una ventaja sustantiva dentro de la lógica de inversión pública, ya que
reduce una de las principales barreras de entrada: la incertidumbre sobre
localización.

Contar con un emplazamiento posible desde etapas tempranas permite ordenar mejor
los estudios de factibilidad, anticipar compatibilidades territoriales e iniciar
conversaciones institucionales con mayor claridad. La localización en Fundo Vaitea
presenta valor territorial por su escala y potencial de proyección, permitiendo
concebir la obra no como una pista aislada, sino como parte de una futura
infra\-estructura deportiva de mayor alcance.

Sin perjuicio de lo anterior, la materialización definitiva deberá quedar sujeta
a los procesos de validación jurídica, factibilidad técnica, compatibilidad
normativa, revisión patrimonial y evaluación territorial que correspondan.

\subsection{Referencia preliminar de costos (Resinsa)}

A la fecha de elaboración del presente informe, se cuenta con una
\textbf{referencia preliminar de costos} asociada a la empresa \textbf{Resinsa},
vinculada a Werner Schmidt Tuki, antecedente que constituye una base inicial
para la estimación temprana del proyecto en etapa de perfil.

Este antecedente debe entenderse únicamente como \textbf{referencia presupuestaria
preeliminar no definitiva}. El costo definitivo deberá establecerse en etapas
posteriores mediante diseño, especificaciones técnicas, cubicaciones y los
procesos de cotización o licitación que correspondan conforme a la normativa
aplicable. Para efectos de postulación, se recomienda trabajar con dos escenarios:

\begin{enumerate}[leftmargin=*, itemsep=3pt]
  \item \textbf{Escenario base:} construcción de pista y operación esencial del recinto
  \item \textbf{Escenario ampliado:} pista con componentes complementarios que
        fortalezcan la visión de Centro Atlético del Pacífico Sur
\end{enumerate}

Esta distinción facilita la priorización de una primera etapa financiable sin
renunciar a la visión estratégica de largo plazo.

% ============================================================
\section{Impacto Ambiental}
% ============================================================

\subsection{Contexto ecosistémico insular}

Rapa Nui es una isla oceánica de origen volcánico, con una ecosistema terrestre
frágil y altamente endémico. Su condición de territorio insular remoto la hace
especialmente sensible a intervenciones de infraestructura que en el continente
tendrían impactos menores.

Los principales factores ecosistémicos a considerar:

\begin{itemize}[leftmargin=*, itemsep=3pt]
  \item Flora endémica con alto grado de vulnerabilidad (incluyendo el \textit{toromiro}
        en peligro de extinción)
  \item Ausencia de rellenos sanitarios continuos: los residuos de construcción
        tienen opciones de disposición limitadas en la isla
  \item Sistemas de agua dulce sensibles (acuífero costero)
  \item Patrimonio arqueológico disperso: cualquier excavación requiere supervisión
        del Consejo de Monumentos Nacionales
\end{itemize}

\subsection{Impactos en fase de construcción}

\begin{table}[h!]
\centering
\begin{tabularx}{\linewidth}{lXX}
\toprule
\textbf{Impacto} & \textbf{Magnitud estimada} & \textbf{Medida de mitigación} \\
\midrule
Movimiento de suelo & Moderado (pista de 400m) & Programa de manejo de tierra y revegetación local \\
Residuos de construcción & Bajo-Moderado & Protocolo de disposición específico para isla; retorno al continente de materiales no reciclables \\
Emisiones de polvo & Bajo (temporal) & Humectación de frentes de trabajo \\
Ruido & Bajo (temporal) & Horarios de obras restringidos \\
Patrimonio arqueológico & Potencialmente Alto & Supervisión del CMN antes y durante excavaciones \\
\bottomrule
\end{tabularx}
\caption{Impactos ambientales en fase de construcción y medidas de mitigación}
\end{table}

\subsection{Huella logística: el costo ambiental de la insularidad}

Un factor ambiental específico de Rapa Nui que documentos estándar omiten: la
\textbf{huella de carbono del transporte de materiales}. Todo material de construcción
---cemento, acero, superficie sintética--- debe llegar a la isla por vía marítima
(barco mensual) o aérea. Este costo logístico ambiental debe considerarse en la
evaluación intuitiva y justifica la preferencia por:

\begin{itemize}[leftmargin=*, itemsep=3pt]
  \item Materiales de alta durabilidad para reducir frecuencia de recambio
  \item Diseño de mínimo mantenimiento complejo
  \item Priorización de mano de obra local para reducir transporte de trabajadores
\end{itemize}

\subsection{Impacto en fase de operación}

Una vez construida, la pista genera impactos ambientales positivos y negativos de
baja magnitud:

\textbf{Positivos:} sustitución de actividades recreativas motorizadas (reconocidas
en la investigación como fuente de lesiones y riesgo), creación de espacio verde
complementario, mejoramiento del paisaje en área urbana de Hanga Roa.

\textbf{Negativos:} aumento de tráfico peatonal y vehicular en el área, consumo
energético de iluminación deportiva. Medidas: iluminación LED de bajo consumo,
programación de horarios de uso.

% ============================================================
\section{Participación Comunitaria y Gobernanza}
\label{sec:gobernanza}
% ============================================================

\subsection{El principio de legitimidad comunitaria}

La investigación cualitativa demostró que el \textbf{primer y más crítico mecanismo}
para el éxito de cualquier iniciativa deportiva en Rapa Nui es la \textbf{legitimidad
local}: que la iniciativa tenga origen y respaldo comunitario explícito, no que sea
percibida como ``de afuera'' \citep{ortega2026}.

Esto no es solo un elemento deseable: es un factor sin el cual las iniciativas
fracasan, independientemente de sus recursos. La investigación documentó el colapso
de iniciativas bien financiadas por falta de respaldo comunitario.

\textbf{Indicador propuesto:} porcentaje de familias y participantes que reconocen
el proyecto como ``nuestro'' versus ``de afuera'', medido en consulta comunitaria
previa al inicio de obras.

\subsection{Actores clave del territorio}

Los siguientes actores deben participar en el proceso de validación y diseño del
proyecto:

\begin{table}[h!]
\centering
\begin{tabularx}{\linewidth}{lX}
\toprule
\textbf{Actor} & \textbf{Rol en el proceso} \\
\midrule
Municipalidad de Isla de Pascua & Coordinación general, facilitador institucional \\
Consejo de Ancianos Rapa Nui (Te Mau Hatu'o) & Validación cultural y territorial, voz de la comunidad histórica \\
Organizaciones deportivas locales & Diagnóstico de necesidades, especificaciones técnicas \\
Comunidades de Hanga Roa & Diagnóstico participativo, usos del espacio \\
Escuelas y liceo local & Programas de uso escolar, formación deportiva \\
Hospital Hanga Roa & Integración salud-deporte, apoyo multidisciplinario \\
IND / Ministerio del Deporte & Marco normativo, financiamiento complementario \\
Gobernación Marítima & Coordinación institucional estatal \\
\bottomrule
\end{tabularx}
\caption{Actores clave y roles en el proceso participativo}
\end{table}

\subsection{Proceso participativo propuesto}

Siguiendo el \textbf{Principio de Legitimidad Primero} derivado de la investigación
\citep{ortega2026}, se propone el siguiente proceso:

\begin{enumerate}[leftmargin=*, itemsep=5pt]
  \item \textbf{Diagnóstico participativo} (antes de diseñar): reuniones con
        dirigentes, familias y jóvenes para validar necesidades y espacios de
        uso. No solo con autoridades.
  \item \textbf{Vocerías reconocidas}: identificar líderes comunitarios que lideran
        el proyecto desde dentro, no solo técnicos externos.
  \item \textbf{Mesa de gobernanza} (durante ejecución): comité con representación
        de clubes deportivos, escuela, salud y municipio, con metas trimestrales
        públicas.
  \item \textbf{Devoluciones pautadas}: reportes a la comunidad del avance
        (mínimo trimestrales), con espacio de corrección antes de decisiones
        mayores.
  \item \textbf{Protocolo de liderazgo femenino}: asegurar participación mínima
        del 40\% de mujeres en todas las instancias, con mentoría y visibilización
        de referentes femeninas locales.
\end{enumerate}

% ============================================================
\section{Barreras y Mecanismos de Traducción}
% ============================================================

\subsection{La teoría ``Del potencial al proyecto''}

La investigación cualitativa que sustenta este informe desarrolló una \textbf{teoría
sustantiva} emergente de los datos: ``\textbf{Del potencial al proyecto}''. Esta
teoría explica que el éxito o fracaso del desarrollo deportivo en Rapa Nui no
depende principalmente de la motivación individual de los atletas, sino de la
\textbf{capacidad colectiva de la comunidad} para activar cinco mecanismos de
traducción bajo condiciones estructurales adversas \citep{ortega2026}.

Aplicada al proyecto de infraestructura, esta teoría ofrece una \textbf{teoría de
cambio} concreta:

\begin{conceptoclave}
  La pista de atletismo no genera impacto social por sí sola. Lo genera en la medida
  en que \textit{habilita} los cinco mecanismos colectivos: legitimidad, profesionalización,
  redes, cuidado integral y gobernanza. Sin estos mecanismos, la infraestructura
  permanecerá subutilizada.
\end{conceptoclave}

\subsection{Los cinco mecanismos de traducción aplicados al proyecto}

\begin{enumerate}[leftmargin=*, itemsep=8pt]
  \item \textbf{Legitimidad local}: La pista debe ser co-diseñada con la comunidad.
    El proceso participativo propuesto en la sección anterior es su condición de
    posibilidad. Sin legitimidad, la infraestructura será percibida como ``de
    afuera'' y su uso será intermitente.

  \item \textbf{Profesionalización gradual}: La pista debe acompañarse de un
    programa de formación de gestores locales ---no sólo de atletas. Transición
    de talleres informales a rutas con roles definidos, estándares y certificación.

  \item \textbf{Redes y alianzas externas con acompañamiento}: La pista reduce
    los costos de entrenamiento local, pero debe acompañarse de convenios con
    federaciones nacionales que reconozcan las limitaciones de la insularidad.
    Negociación de excepciones fundamentadas al régimen colonial-burocrático.

  \item \textbf{Cuidado integral del atleta}: La pista debe ser parte de un
    ecosistema que incluya prevención de lesiones, apoyo mental básico y
    orientación nutricional. El \textbf{Equipo Mínimo Integral} propuesto por
    la investigación: entrenador/a certificado/a, apoyo en salud mental,
    kinesiólogo/a, orientación nutricional.

  \item \textbf{Gobernanza coordinada}: Mesa interactor (clubes-escuela-salud-municipio)
    con metas anuales públicas, criterios de selección de atletas transparentes
    y seguimiento de cohortes ---no solo de ``estrellas''.
\end{enumerate}

% ============================================================
\section{Análisis de Alternativas}
% ============================================================

\subsection{Alternativa 1: Mantener situación actual}

Consiste en continuar utilizando infraestructura no especializada (calles,
espacio frente al estadio, otras superficies). Esta alternativa presenta el
menor costo inmediato, pero perpetúa la brecha documentada, limita el
desarrollo deportivo y mantiene el déficit de calidad en la oferta pública.
Su impacto sobre las trayectorias deportivas es nulo, y agrava la deserción
documentada en la investigación cualitativa \citep{ortega2026}.

\subsection{Alternativa 2: Mejoramiento parcial de espacios existentes}

Implica adaptar o mejorar espacios ya disponibles (estadio, multicanchas).
Puede tener menor costo inicial, pero usualmente produce rendimientos limitados,
usos incompatibles y una vida útil funcional menor en términos atléticos. No
resuelve el problema estructural de falta de pista certificada.

\subsection{Alternativa 3: Construcción de pista de atletismo con enfoque multipropósito}

Es la alternativa recomendada. Permite resolver de manera estructural el problema,
aumentar la capacidad comunal, generar mayor impacto social y construir una obra
escalable hacia el Centro Atlético del Pacífico Sur.

\subsection{Síntesis comparativa}

\begin{table}[h!]
\centering
\caption{Comparación cualitativa de alternativas}
\begin{tabularx}{\linewidth}{
  >{\RaggedRight\arraybackslash}p{3.5cm}
  >{\centering\arraybackslash}X
  >{\centering\arraybackslash}X
  >{\centering\arraybackslash}X}
\toprule
\textbf{Criterio} & \textbf{Situación actual} & \textbf{Mejoramiento parcial} & \textbf{Pista multipropósito} \\
\midrule
Resolución del problema & Baja & Media & Alta \\
Impacto social & Bajo & Medio & Alto \\
Escalabilidad & Nula & Baja & Alta \\
Visibilidad política & Baja & Media & Alta \\
Potencial deportivo & Bajo & Medio & Alto \\
Aporte al SNI & Bajo & Medio & Alto \\
Compatibilidad con investigación & No & Parcial & Total \\
\bottomrule
\end{tabularx}
\end{table}

% ============================================================
\section{Hoja de Ruta para Formulación y Postulación}
% ============================================================

\subsection{Fase 1: Validación política e institucional}

\begin{enumerate}[leftmargin=*, itemsep=3pt]
  \item Aprobación conceptual municipal
  \item Definición del relato estratégico oficial (narrativa para SNI)
  \item Priorización del proyecto en cartera comunal/regional
\end{enumerate}

\subsection{Fase 2: Validación territorial y comunitaria}

\begin{enumerate}[leftmargin=*, itemsep=3pt]
  \item Revisión jurídica del terreno en Fundo Vaitea
  \item Diagnóstico participativo con actores comunitarios (ver Sección~\ref{sec:gobernanza})
  \item Evaluación inicial de compatibilidades territoriales y patrimoniales
\end{enumerate}

\subsection{Fase 3: Estudios habilitantes}

\begin{enumerate}[leftmargin=*, itemsep=3pt]
  \item Topografía del emplazamiento propuesto
  \item Mecánica de suelos
  \item Inspección arqueológica (Consejo de Monumentos Nacionales)
  \item Perfil técnico-económico con costos de ciclo de vida
  \item Presupuesto detallado con sobrecosto insular
\end{enumerate}

\subsection{Fase 4: Formulación y postulación}

\begin{enumerate}[leftmargin=*, itemsep=3pt]
  \item Ingreso de iniciativa al Sistema Nacional de Inversiones (SNI)
  \item Formulación bajo metodología sectorial deportes \citep{SNIdeportes}
  \item Validación presupuestaria y búsqueda de financiamiento
\end{enumerate}

\subsection{Fase 5: Diseño, ejecución y activación comunitaria}

\begin{enumerate}[leftmargin=*, itemsep=3pt]
  \item Licitación de diseño con participación comunitaria
  \item Construcción con mano de obra local priorizada
  \item Reglamento y plan de operación validado por la comunidad
  \item Calendario de uso desde la inauguración
\end{enumerate}

% ============================================================
\section{Argumentario Institucional para Autoridades}
% ============================================================

Este proyecto tiene una narrativa sólida y diferenciada según la audiencia
institucional a la que se presente.

\begin{description}[leftmargin=3.5cm, style=nextline, itemsep=6pt]

  \item[Municipalidad] \textit{Este proyecto entrega a la comuna una infraestructura
    visible, útil y transversal, con capacidad de impacto real sobre niñez, juventud,
    salud y vida comunitaria. Es un legado tangible de inversión pública.}

  \item[Gob. Regional] \textit{La iniciativa reduce brechas territoriales en un territorio
    extremo, combina equidad con alto valor simbólico y puede transformarse en una
    obra emblemática de descentralización efectiva.}

  \item[IND / Deporte] \textit{La pista crea base física para programas formativos, deporte
    escolar, desarrollo atlético y articulación con otras disciplinas. Es inversión con
    retorno deportivo medible.}

  \item[Patrimonio] \textit{El proyecto puede ser una obra pública contemporánea, sobria
    y respetuosa, diseñada con participación local y compatible con la identidad
    territorial y el paisaje cultural de Rapa Nui.}

  \item[Ciudadanía] \textit{La pista es de todos. Es una inversión que mejora la calidad
    de vida de niños, jóvenes, familias y adultos mayores de Rapa Nui, con respeto
    por la cultura y el territorio.}

\end{description}

% ============================================================
\section{Conclusión y Justificación de la Inversión}
% ============================================================

La construcción de una pista de atletismo en Isla de Pascua representa una inversión
pública estratégica que contribuiría al desarrollo social, deportivo, cultural y
comunitario de un territorio con características únicas en el sistema de inversión
pública chileno.

\subsection{Síntesis de impactos}

\begin{table}[h!]
\centering
\begin{tabularx}{\linewidth}{lXl}
\toprule
\textbf{Dimensión} & \textbf{Impacto esperado} & \textbf{Magnitud} \\
\midrule
Salud pública & Reducción sedentarismo, prevención enfermedades crónicas, salud mental & Alta \\
Formación deportiva & Trayectorias sostenidas, reducción deserción & Alta \\
Cohesión comunitaria & Espacio de encuentro, capital social & Alta \\
Juventud & Reducción conductas riesgo, disciplina, oportunidades & Alta \\
Identidad cultural & Práctica deportiva como acto identitario rapanui & Alta \\
Género & Acceso equitativo con diseño participativo & Media-Alta \\
Ambiental (construcción) & Bajo, controlable con mitigaciones específicas & Baja \\
Ambiental (operación) & Mínimo, positivo neto & Muy baja \\
\bottomrule
\end{tabularx}
\caption{Resumen de impactos por dimensión}
\end{table}

\subsection{Condición crítica para el éxito}

La evidencia empírica de esta investigación es clara: ninguna infraestructura
deportiva en Rapa Nui puede maximizar su impacto sin un \textbf{proceso participativo
genuino} que preceda y acompañe su instalación. La pista como objeto material es
condición necesaria pero no suficiente.

La recomendación es avanzar en la formulación del proyecto dentro del
\textbf{Sistema Nacional de Inversiones}, incorporando:

\begin{enumerate}[leftmargin=*, itemsep=4pt]
  \item Estudios técnicos y arquitectónicos en diálogo con la comunidad
  \item Proceso participativo con los actores identificados en la Sección 7
  \item Validación comunitaria del diseño antes de licitación
  \item Diseño de un programa de uso y gobernanza post-construcción (no solo
        construcción del recinto)
  \item Mecanismo de monitoreo y evaluación con indicadores acordados con
        la comunidad (retención, participación femenina, continuidad escolar,
        incidencia de lesiones)
\end{enumerate}

\begin{center}
  \textit{``El problema no es `tener o no tener talento', sino cómo la comunidad
  convierte ese talento en un proyecto sostenible mientras negocia estructuras
  que no fueron hechas para ella.''}\\[6pt]
  \small--- Síntesis de la investigación cualitativa, Ortega \& Nervi (2026)
\end{center}

Considerando el aislamiento geográfico de la isla, la disponibilidad limitada de
infraestructura deportiva, la historia colonial documentada y la relevancia del
deporte como herramienta de desarrollo social y cultural, el proyecto presenta una
\textbf{alta pertinencia social, territorial y cultural dentro del marco de inversión
pública en Chile}.

% ============================================================
% ANEXOS
% ============================================================
\newpage
\section*{Anexo A: Propuesta de cartera inicial de usos del recinto}
\addcontentsline{toc}{section}{Anexo A: Cartera inicial de usos del recinto}

\begin{table}[h!]
\centering
\caption{Propuesta de programación por línea de uso}
\begin{tabularx}{\linewidth}{>{\bfseries}p{3.5cm}X}
\toprule
\textbf{Línea programática} & \textbf{Uso propuesto} \\
\midrule
Deporte escolar & Pruebas atléticas, talleres y encuentros inter escolares. \\
Formación deportiva & Escuela municipal de atletismo por categorías etarias. \\
Salud comunitaria & Caminatas guiadas, acondicionamiento físico, programas para adultos. \\
Juventud & Talleres vespertinos de velocidad, resistencia y fuerza básica. \\
Comunidad & Corridas familiares, jornadas de actividad física, festivales deportivos. \\
Turismo deportivo & Clínicas, campamentos y programas de intercambio de escala compatible. \\
\bottomrule
\end{tabularx}
\end{table}

\section*{Anexo B: Información mínima para presentación institucional}
\addcontentsline{toc}{section}{Anexo B: Información mínima para presentación institucional}

Antes de entregar este documento como versión final ante organismos de inversión
pública, se recomienda completar los siguientes elementos con información local
validada:

\begin{itemize}[leftmargin=*, itemsep=3pt]
  \item Terreno propuesto y justificación de localización (Fundo Vaitea, verificación jurídica)
  \item Croquis o plano de emplazamiento
  \item Fotografías del área actual
  \item Carta de apoyo municipal o institucional firmada
  \item Estimación preliminar de costos (Resinsa, actualizada)
  \item Número de establecimientos educacionales beneficiarios
  \item Programas deportivos activos en la comuna (diagnóstico actualizado 2025--2026)
  \item Diagnóstico breve de uso esperado por segmento etario
  \item Ruta de permisos y compatibilidades patrimoniales (CMN, CONAF)
\end{itemize}

\section*{Anexo C: Antecedentes preliminares de formulación}
\addcontentsline{toc}{section}{Anexo C: Antecedentes preliminares de formulación}

Para efectos de la presente versión del informe, se consideran como antecedentes
preliminares de formulación los siguientes elementos de trabajo que deben ser
formalizados en etapas posteriores:

\begin{itemize}[leftmargin=*, itemsep=3pt]
  \item Identificación de un emplazamiento preliminar en el sector Fundo Vaitea
  \item Existencia de terrenos donados para fines asociados al proyecto
  \item Referencia preliminar de costos asociada a la empresa Resinsa, vinculada
        a Werner Schmidt Tuki, utilizada únicamente como base inicial de
        dimensionamiento de inversión
\end{itemize}

Estos antecedentes deberán ser formalizados, validados y complementados en las
etapas posteriores de prefactibilidad, diseño y formulación oficial de la
iniciativa ante el SNI.

% ============================================================
% GLOSARIO
% ============================================================
\newpage
\section*{Glosario de Términos Rapanui}
\addcontentsline{toc}{section}{Glosario de términos rapanui}

\begin{description}[leftmargin=3cm, style=nextline, itemsep=6pt]
  \item[\textit{Hanai}] Práctica de crianza delegada donde tíos, abuelos u otros
    miembros de la familia extendida asumen responsabilidad de crianza. La relación
    filial puede ser instaurada por el acto del cuidado y la alimentación, no por
    vínculo biológico.
  \item[\textit{Vā}] Espacio relacional entre personas que debe ser cuidado y
    nutrido. No es un espacio físico sino relacional; su cuidado es condición de
    la vida comunitaria.
  \item[\textit{Talanoa}] Metodología conversacional polinésica que prioriza la
    reciprocidad y el flujo relacional por sobre la rigidez de una guía de preguntas
    fija. Fue adoptada como principio metodológico de la investigación.
  \item[\textit{Koros}] Abuelos, ancianos conocedores. Transmisores de saberes
    técnicos y culturales, incluyendo conocimientos deportivos tradicionales.
  \item[\textit{Tapati Rapa Nui}] Festival cultural anual de la isla que incluye
    componentes deportivos, artísticos y de destreza física. Expresión viva de la
    identidad rapanui contemporánea.
  \item[\textit{Haka Pei}] Prueba de valentía del Tapati: competidores se deslizan
    por la ladera de un cerro sobre troncos de plátano, a velocidades de hasta 80 km/h.
  \item[\textit{Tangata Manu}] Antiguo rito del Hombre Pájaro, última gran expresión
    de la cultura pre-cristiana en Rapa Nui (1866). Combinaba destreza física,
    valentía y favor divino.
  \item[\textit{Te Mau Hatu'o Rapa Nui}] Consejo de Ancianos de Rapa Nui,
    reconstituido en 1983. Voz histórica de la comunidad en defensa de sus
    derechos territoriales.
  \item[\textit{Papa Ra'au}] Medicina tradicional rapanui, en tensión y diálogo
    con la medicina occidental en el Hospital Hanga Roa.
\end{description}

% ============================================================
% BIBLIOGRAFÍA
% ============================================================
\newpage
\bibliography{referencias}

\end{document}
